% Part: first-order-logic
% Chapter: syntax-and-semantics
% Section: assignments

\documentclass[../../../include/open-logic-section]{subfiles}

\begin{document}

\olfileid{fol}{syn}{ass}

\olsection{Pangaji Variabel}

\begin{explain}
Pangaji !!{variable}~$s$ menehi nilai marang \emph{saben}
variabel---lan cacah variabel iku tanpa wates. Mesthi wae,
syarat iki ora perlu. Kita njaluk supaya pangaji !!{variable}
menehi nilai marang kabeh !!{variable}s mung amarga cara iki
nggampangake prakara. Nilai term~$t$, lan apa !!a{formula}~$!A$
dicukupi ing !!a{structure} relatif marang~$s$, mung gumantung
marang nilai kang diwenehake~$s$ marang !!{variable}s ing~$t$
lan !!{variable}s bebas ing~$!A$. Iki isi saka rong proposisi
sabanjure. Supaya teges "gumantung marang" cetha, kita nuduhake
yen rong pangaji variabel kang padha ing kabeh variabel
ing~$t$ menehi nilai kang padha, lan yen rong pangaji variabel
padha ing kabeh variabel bebas saka~$!A$, mula $!A$
dicukupi relatif marang pangaji kapisan yen lan mung yen
dicukupi relatif marang pangaji kapindho.
\end{explain}

\begin{prop}\ollabel{prop:valindep}
Yen !!{variable}s ing term~$t$ kalebu $x_1$, \dots,~$x_n$, lan
$s_1(x_i) = s_2(x_i)$ kanggo $i = 1$, \dots,~$n$, mula $\Value{t}{M}[s_1]
= \Value{t}{M}[s_2]$.
\end{prop}

\begin{proof}
Kanthi induksi miturut keruwetan~$t$. Kanggo kasus dhasar, $t$ bisa
dadi !!a{constant} utawa salah siji variabel~$x_1$, \dots,~$x_n$.
Yen $t = c$, mula $\Value{t}{M}[s_1] = \Assign{c}{M} =
\Value{t}{M}[s_2]$. Yen $t = x_i$, $s_1(x_i) = s_2(x_i)$ miturut
hipotesis proposisi, mula $\Value{t}{M}[s_1] = s_1(x_i) =
s_2(x_i) = \Value{t}{M}[s_2]$.

Kanggo langkah induksi, upamakna $t = \Atom{f}{t_1, \dots, t_k}$
lan pratelan iki bener kanggo $t_1$, \dots, $t_k$. Mula
\begin{align*}
  \Value{t}{M}[s_1] & = \Value{\Atom{f}{t_1, \dots, t_k}}{M}[s_1] \\
  & = \Assign{f}{M}(\Value{t_1}{M}[s_1], \dots, \Value{t_k}{M}[s_1]).
\intertext{Kanggo $j = 1$, \dots,~$k$, !!{variable}s ing~$t_j$
    kalebu $x_1$, \dots,~$x_n$. Miturut hipotesis induksi,
    $\Value{t_j}{M}[s_1] = \Value{t_j}{M}[s_2]$. Dadi,}
\Value{t}{M}[s_1] & = \Value{\Atom{f}{t_1, \dots, t_k}}{M}[s_1] \\
  & = \Assign{f}{M}(\Value{t_1}{M}[s_1], \dots, \Value{t_k}{M}[s_1]) \\
  & = \Assign{f}{M}(\Value{t_1}{M}[s_2], \dots, \Value{t_k}{M}[s_2]) \\
  & = \Value{\Atom{f}{t_1, \dots, t_k}}{M}[s_2] = \Value{t}{M}[s_2].
\end{align*}
\end{proof}

\begin{prop}\ollabel{prop:satindep}
Yen !!{variable}s bebas ing $!A$ kalebu $x_1$, \dots,~$x_n$, lan
$s_1(x_i) = s_2(x_i)$ kanggo $i = 1$, \dots,~$n$, mula $\Sat{M}{!A}[s_1]$
yen lan mung yen $\Sat{M}{!A}[s_2]$.
\end{prop}

\begin{proof}
Kita nggunakake induksi miturut keruwetan $!A$. Ing kasus dhasar,
nalika $!A$ atom, $!A$ bisa awujud:
\iftag{prvTrue}{$\ltrue$,}{}
\iftag{prvFalse}{$\lfalse$,}{}
$\Atom{R}{t_1, \dots, t_k}$ kanggo predikat $k$-panggonan $R$
lan term $t_1$, \dots,~$t_k$, utawa $\eq[t_1][t_2]$ kanggo term
$t_1$ lan~$t_2$. Ing rong kasus pungkasan, kita mung nuduhake
arah maju saka !!{biconditional}, amarga bukti arah walike simetris.

\begin{enumerate}
\tagitem{prvTrue}{%
  \indcase{!A}{\ltrue}{loro-lorone $\Sat{M}{\indfrm}[s_1]$ lan
    $\Sat{M}{\indfrm}[s_2]$.}}{}

\tagitem{prvFalse}{%
  \indcase{!A}{\lfalse}{loro-lorone $\Sat/{M}{!A}[s_1]$ lan
    $\Sat/{M}{!A}[s_2]$.}}{}

\item
  \indcase{!A}{\Atom{R}{t_1, \ldots, t_k}}{upamakna
    $\Sat{M}{\indfrm}[s_1]$. Mula
    \[
    \langle \Value{t_1}{M}[s_1], \ldots, \Value{t_k}{M}[s_1] \rangle
    \in \Assign{R}{M}.
    \]
    Kanggo $i = 1$, \dots,~$k$, $\Value{t_i}{M}[s_1] =
    \Value{t_i}{M}[s_2]$ miturut \olref{prop:valindep}. Mula uga
    $\langle \Value{t_1}{M}[s_2], \ldots, \Value{t_k}{M}[s_2] \rangle
    \in \Assign{R}{M}$, saengga $\Sat{M}{\indfrm}[s_2]$.
    Cathetan koreksi sumber OLPL-079: entri kapisan ing
    tupel panutup sumber migunakake indeks umum, dudu indeks
    kapisan; indeks wiwitan kang cocog dibalekake.}
\item
  \indcase{!A}{\eq[t_1][t_2]}{upamakna $\Sat{M}{\indfrm}[s_1]$.
    Mula $\Value{t_1}{M}[s_1] = \Value{t_2}{M}[s_1]$. Dadi,
    \begin{align*}
      \Value{t_1}{M}[s_2] & = \Value{t_1}{M}[s_1]
      & \text{(miturut \olref{prop:valindep})} \\
      & = \Value{t_2}{M}[s_1]
      & \text{(amarga $\Sat{M}{\eq[t_1][t_2]}[s_1]$)}\\
      &= \Value{t_2}{M}[s_2]
      & \text{(miturut \olref{prop:valindep}),}
    \end{align*}
    mula $\Sat{M}{\eq[t_1][t_2]}[s_2]$.}
\end{enumerate}

Saiki upamakna $\Sat{M}{!B}[s_1]$ yen lan mung yen
$\Sat{M}{!B}[s_2]$ kanggo kabeh !!{formula}s~$!B$
kang kurang rumit tinimbang~$!A$. Langkah induksi dipérang
miturut operator utama~$!A$. Ing saben kasus, kita mung
nuduhake arah maju saka !!{biconditional}; bukti arah
walike simetris. Ing kabeh kasus kajaba kuantor, kita
ngetrapake hipotesis induksi marang sub-!!{formula}s~$!B$
saka~$!A$. Variabel bebas ing~$!B$ kalebu variabel bebas
ing~$!A$. Mula, yen $s_1$ lan $s_2$ padha ing variabel bebas
ing~$!A$, loro pangaji iku uga padha ing variabel bebas
ing~$!B$, saengga hipotesis induksi bisa ditrapake marang~$!B$.

\begin{enumerate}
\tagitem{defNot}{}{%
  \iftag{probNot}{%
    \indcase!{!A}{\lnot !B}{}}{%
    \indcase{!A}{\lnot !B}{yen $\Sat{M}{\indfrm}[s_1]$, mula
      $\Sat/{M}{!B}[s_1]$, dadi miturut hipotesis induksi,
      $\Sat/{M}{!B}[s_2]$, saengga $\Sat{M}{\indfrm}[s_2]$.}}}

\tagitem{defAnd}{}{%
  \iftag{probAnd}{%
    \indcase!{!A}{!B \land !C}{}}{%
    \indcase{!A}{!B \land !C}{yen $\Sat{M}{\indfrm}[s_1]$, mula
      $\Sat{M}{!B}[s_1]$ lan $\Sat{M}{!C}[s_1]$, mula miturut
      hipotesis induksi, $\Sat{M}{!B}[s_2]$ lan
      $\Sat{M}{!C}[s_2]$. Dadi $\Sat{M}{\indfrm}[s_2]$.}}}

\tagitem{defOr}{}{%
  \iftag{probOr}{%
    \indcase!{!A}{!B \lor !C}{}}{%
    \indcase{!A}{!B \lor !C}{yen $\Sat{M}{\indfrm}[s_1]$, mula
      $\Sat{M}{!B}[s_1]$ utawa $\Sat{M}{!C}[s_1]$. Miturut
      hipotesis induksi, $\Sat{M}{!B}[s_2]$ utawa
      $\Sat{M}{!C}[s_2]$, saengga $\Sat{M}{\indfrm}[s_2]$.}}}

\tagitem{defIf}{}{%
  \iftag{probIf}{%
    \indcase!{!A}{!B \lif !C}{}}{%
    \indcase{!A}{!B \lif !C}{yen $\Sat{M}{\indfrm}[s_1]$, mula
    $\Sat/{M}{!B}[s_1]$ utawa $\Sat{M}{!C}[s_1]$. Miturut
    hipotesis induksi, $\Sat/{M}{!B}[s_2]$ utawa
    $\Sat{M}{!C}[s_2]$, saengga $\Sat{M}{\indfrm}[s_2]$.}}}

\tagitem{defIff}{}{%
  \iftag{probIff}{%
    \indcase!{!A}{!B \liff !C}{}}{%
    \indcase{!A}{!B \liff !C}{yen $\Sat{M}{\indfrm}[s_1]$, mula
      salah siji $\Sat{M}{!B}[s_1]$ lan $\Sat{M}{!C}[s_1]$,
      utawa $\Sat/{M}{!B}[s_1]$ lan $\Sat/{M}{!C}[s_1]$.
      Miturut hipotesis induksi, salah siji
      $\Sat{M}{!B}[s_2]$ lan $\Sat{M}{!C}[s_2]$ utawa
      $\Sat/{M}{!B}[s_2]$ lan $\Sat/{M}{!C}[s_2]$.
      Ing salah siji kasus kasebut, $\Sat{M}{\indfrm}[s_2]$.}}}

\tagitem{defEx}{}{%
  \iftag{probEx}{%
    \indcase!{!A}{\lexists[x][!B]}{}}{%
    \indcase{!A}{\lexists[x][!B]}{yen $\Sat{M}{\indfrm}[s_1]$,
      ana $m \in \Domain{M}$ kang ndadekake
      $\Sat{M}{!B}[\Subst{s_1}{m}{x}]$.
      Upamakna $s_1' = \Subst{s_1}{m}{x}$ lan $s_2' =
      \Subst{s_2}{m}{x}$. Variabel bebas ing~$!B$ kalebu $x_1$,
      \dots, $x_n$, lan $x$. $s_1'(x_i) = s_2'(x_i)$, amarga
      $s_1'$ lan $s_2'$ iku varian-$x$ saka $s_1$ lan~$s_2$
      miturut urutane, lan miturut hipotesis
      $s_1(x_i) = s_2(x_i)$. $s_1'(x) = s_2'(x) = m$
      miturut cara kita netepake $s_1'$ lan~$s_2'$.
      Mula hipotesis induksi bisa ditrapake marang $!B$ lan
      $s_1'$, $s_2'$, saengga $\Sat{M}{!B}[s_2']$.
      Dadi, amarga $s_2' = \Subst{s_2}{m}{x}$,
      ana $m \in \Domain{M}$ kang ndadekake
      $\Sat{M}{!B}[\Subst{s_2}{m}{x}]$, saengga
      $\Sat{M}{\indfrm}[s_2]$.}}}

\tagitem{defAll}{}{%
  \iftag{probAll}{%
    \indcase!{!A}{\lforall[x][!B]}{}}{%
    \indcase{!A}{\lforall[x][!B]}{yen $\Sat{M}{\indfrm}[s_1]$,
      mula kanggo saben $m \in \Domain{M}$,
      $\Sat{M}{!B}[\Subst{s_1}{m}{x}]$. Kita arep nuduhake yen
      uga kanggo saben $m \in \Domain{M}$,
      $\Sat{M}{!B}[\Subst{s_2}{m}{x}]$. Dadi pilih
      $m \in \Domain{M}$ sembarang, lan gatekna
      $s_1' = \Subst{s_1}{m}{x}$ lan $s_2' =
      \Subst{s_2}{m}{x}$. Cathetan koreksi sumber OLPL-080:
      rong pangaji turunan ing sumber nganggo pangaji dhasar
      kang ora ditetepake; saben turunan saiki nganggo pangaji
      dhasare dhewe. Kita nduweni $\Sat{M}{!B}[s_1']$.
      Variabel bebas ing~$!B$ kalebu $x_1$, \dots, $x_n$, lan $x$.
      $s_1'(x_i) = s_2'(x_i)$, amarga $s_1'$ lan $s_2'$ iku
      varian-$x$ saka $s_1$ lan~$s_2$ miturut urutane, lan miturut
      hipotesis $s_1(x_i) = s_2(x_i)$. $s_1'(x) = s_2'(x) = m$
      miturut cara kita netepake $s_1'$ lan~$s_2'$. Mula hipotesis
      induksi bisa ditrapake marang~$!B$ lan~$s_1'$, $s_2'$,
      lan kita nduweni $\Sat{M}{!B}[s_2']$. Iki bener
      kanggo saben $m \in \Domain{M}$, yaiku
      $\Sat{M}{!B}[\Subst{s_2}{m}{x}]$ kanggo kabeh~$m \in
      \Domain{M}$, saengga $\Sat{M}{\indfrm}[s_2]$.}}}
\end{enumerate}
Kanthi induksi, kita oleh $\Sat{M}{!A}[s_1]$ yen lan mung yen
$\Sat{M}{!A}[s_2]$ kapan wae !!{variable}s bebas ing $!A$
kalebu $x_1$, \dots, $x_n$ lan $s_1(x_i)=s_2(x_i)$
kanggo $i = 1$, \dots,~$n$.
\end{proof}

\begin{probtag}{probNot,probOr,probAnd,probIf,probIff,probEx,probAll}
Rampungna bukti \olref[fol][syn][ass]{prop:satindep}.
\end{probtag}

\begin{explain}
!!^{sentence}s ora nduweni variabel bebas, mula sembarang
rong pangaji variabel menehi nilai kang padha marang kabeh
variabel bebas (kang cacahé nol) saka saben ukara. Dadi
proposisi kang mentas dibuktekake ngandharake yen apa
!!a{sentence} dicukupi ing struktur relatif marang pangaji
variabel pancen ora gumantung marang pangaji kasebut.
Kita bakal nyathet kasunyatan iki. Iki dadi dhasar
definisi sabanjure bab relasi nyukupi !!a{sentence} ing
!!a{structure} tanpa nyebut pangaji variabel.
\end{explain}

\begin{cor}
\ollabel{cor:sat-sentence}
Yen $!A$ iku !!a{sentence} lan $s$ pangaji variabel, mula
$\Sat{M}{!A}[s]$ yen lan mung yen $\Sat{M}{!A}[s']$
kanggo saben pangaji variabel~$s'$.
\end{cor}

\begin{proof}
  Upamakna $s'$ sembarang pangaji variabel. Amarga $!A$ iku
  !!a{sentence}, ukara iku ora nduweni variabel bebas,
  mula saben pangaji variabel~$s'$ kanthi otomatis menehi
  nilai kang padha marang kabeh variabel bebas ing~$!A$
  kaya kang diwenehake~$s$. Dadi syarat
  \olref{prop:satindep} kacukupan, lan kita oleh
  $\Sat{M}{!A}[s]$ yen lan mung yen $\Sat{M}{!A}[s']$.
\end{proof}

\begin{defn}
\ollabel{defn:satisfaction}
Yen $!A$ iku !!a{sentence}, kita kandha yen
!!a{structure}~$\Struct M$ \emph{nyukupi}~$!A$,
$\Sat{M}{!A}$, yen lan mung yen $\Sat{M}{!A}[s]$
kanggo kabeh pangaji variabel~$s$.
\end{defn}

Yen $\Sat{M}{!A}$, kita uga kandha kanthi prasaja yen
\emph{$!A$ bener ing~$\Struct{M}$.} Gagasan relasi nyukupi
lumrahé bisa dijembarake saka !!{sentence}s siji-siji
marang himpunan !!{sentence}s.

\begin{defn}
\ollabel{defn:sat}
Yen $\Gamma$ iku himpunan !!{sentence}s~$\Gamma$, kita
kandha yen !!a{structure}~$\Struct M$ \emph{nyukupi}~$\Gamma$,
$\Sat{M}{\Gamma}$, yen lan mung yen $\Sat{M}{!A}$
kanggo kabeh $!A \in \Gamma$.
\end{defn}

\begin{prop}\ollabel{prop:sentence-sat-true}
  Upamakna $\Struct{M}$ iku !!a{structure}, $!A$ iku
  !!a{sentence}, lan $s$ pangaji variabel.
  $\Sat{M}{!A}$ yen lan mung yen $\Sat{M}{!A}[s]$.
\end{prop}

\begin{proof}
Gladhen.
\end{proof}

\begin{prob}
Buktèkna \olref[fol][syn][ass]{prop:sentence-sat-true}
\end{prob}

\begin{prop}\ollabel{prop:sat-quant}
  Upamakna $!A(x)$ mung ngemot $x$ minangka variabel bebas,
  lan $\Struct M$ iku !!a{structure}. Mula:
  \begin{tagenumerate}{prvEx,prvAll}
    \tagitem{prvEx}{$\Sat{M}{\lexists[x][!A(x)]}$ yen lan mung yen
      $\Sat{M}{!A(x)}[s]$ kanggo paling ora siji pangaji variabel~$s$.}{}
    \tagitem{prvAll}{$\Sat{M}{\lforall[x][!A(x)]}$ yen lan mung yen
      $\Sat{M}{!A(x)}[s]$ kanggo kabeh pangaji variabel~$s$.}{}
\end{tagenumerate}
\end{prop}

\begin{proof}
Gladhen.
\end{proof}

\begin{prob}
Buktèkna \olref[fol][syn][ass]{prop:sat-quant}.
\end{prob}

\begin{prob}
\DeclareRobustCommand{\VDash}{\mathrel{||}\joinrel\Relbar}
Upamakna $\Lang L$ iku basa tanpa !!{function}s. Kanggo
!!{structure}~$\Struct M$, $c$ !!a{constant}, lan
$a \in \Domain M$, tetapna $\Struct M[a/c]$ minangka
!!{structure} kang padha karo~$\Struct M$ kajaba
$\Assign{c}{M[a/c]} = a$. Tetepna $\Struct M \VDash !A$
kanggo !!{sentence}s~$!A$ kaya mangkene:
\begin{enumerate}
\tagitem{prvFalse}{%
  \indcase{!A}{\lfalse}{ora $\Struct{M} \VDash \indfrm$.}}{}

\tagitem{prvTrue}{%
  \indcase{!A}{\ltrue}{$\Struct{M} \VDash \indfrm$.}}{}

\item \indcase{!A}{\Atom{R}{d_1, \dots, d_n}}{$\Struct M \VDash \indfrm$
  yen lan mung yen $\langle \Assign{d_1}{M}, \dots, \Assign{d_n}{M} \rangle \in
  \Assign{R}{M}$.}

\item \indcase{!A}{\eq[d_1][d_2]}{$\Struct M \VDash \indfrm$ yen lan mung yen
  $\Assign{d_1}{M} = \Assign{d_2}{M}$.}

\tagitem{prvNot}{%
  \indcase{!A}{\lnot !B}{$\Struct{M} \VDash \indfrm$ yen lan mung yen
    ora $\Struct{M} \VDash {!B}$.}}{}

\tagitem{prvAnd}{%
  \indcase{!A}{(!B \land !C)}{$\Struct{M} \VDash
    \indfrm$ yen lan mung yen $\Struct{M} \VDash!B$
    lan $\Struct{M} \VDash !C$.}}{}

\tagitem{prvOr}{%
  \indcase{!A}{(!B \lor !C)}{$\Struct{M} \VDash \indfrm$ yen lan mung yen
    $\Struct{M} \VDash !B$ utawa $\Struct{M} \VDash !C$
    (utawa kalorone).}}{}

\tagitem{prvIf}{%
  \indcase{!A}{(!B \lif !C)}{$\Struct{M} \VDash \indfrm$ yen lan mung yen
    ora $\Struct {M} \VDash !B$ utawa $\Struct M \VDash !C$
    (utawa kalorone).}}{}

\tagitem{prvIff}{%
  \indcase{!A}{(!B \liff !C)}{$\Struct{M} \VDash {\indfrm}$ yen lan mung yen
    salah siji $\Struct{M} \VDash {!B}$ lan
    $\Struct{M} \VDash {!C}$ kalorone, utawa
    $\Struct{M} \VDash {!B}$ lan $\Struct{M} \VDash {!C}$
    kalorone ora bener.}}{}

\tagitem{prvAll}{%
  \indcase{!A}{\lforall[x][!B]}{$\Struct{M} \VDash {\indfrm}$
    yen lan mung yen kanggo kabeh $a \in \Domain{M}$,
    $\Struct{M[a/c]} \VDash \Subst{!B}{c}{x}$,
    yen $c$ ora katon ing~$!B$.}}{}

\tagitem{prvEx}{%
  \indcase{!A}{\lexists[x][!B]}{$\Struct{M} \VDash
  {\indfrm}$ yen lan mung yen ana $a \in \Domain M$ kang
  ndadekake $\Struct{M[a/c]} \VDash \Subst{!B}{c}{x}$,
  yen $c$ ora katon ing~$!B$.}}{}
\end{enumerate}
Upamakna $x_1$, \dots, $x_n$ iku kabeh !!{variable}s bebas
ing~$!A$, $c_1$, \dots, $c_n$ simbol konstanta kang ora ana
ing~$!A$, $a_1$, \dots, $a_n \in \Domain M$, lan
$s(x_i) = a_i$.

Tuduhna yen $\Sat{M}{!A}[s]$ yen lan mung yen
$\Struct M[a_1/c_1,\dots,a_n/c_n]
\VDash \Subst{\Subst{!A}{c_1}{x_1}\dots}{c_n}{x_n}$.

(Masalah iki nuduhake yen semantik logika sepisanan bisa
ditetepake tanpa nggunakake pangaji variabel.)
\end{prob}

\begin{prob}
Upamakna $f$ iku simbol fungsi kang ora ana
ing~$!A(x,y)$. Tuduhna yen ana !!a{structure}~$\Struct{M}$
kang ndadekake
$\Sat{M}{\lforall[x][\lexists[y][!A(x,y)]]}$
yen lan mung yen ana~$\Struct M'$ kang ndadekake
$\Sat{M'}{\lforall[x][!A(x,f(x))]}$.

(Masalah iki kasus mligi saka kang diarani Teorema Skolem;
$\lforall[x][!A(x,f(x))]$ diarani \emph{wangun normal Skolem}
saka $\lforall[x][\lexists[y][!A(x,y)]]$.)
\end{prob}

\end{document}
